\documentclass[../../main.tex]{subfiles}
\graphicspath{{../../image/}{../../figure/}}
\begin{document}

\section{全局波型}\label{section:ch14-全局波型}

从某固定初始阶梯函数出发，\cref{section:ch14-真非线性系统的前沿追踪}中所述前沿追踪算法将在一个极大时间区间 $[0,T)$ 上产生分段常值函数 $U$。原则上 $T$ 可能有限，即碰撞数当 $t\uparrow T$ 时无界增长；任务是证明此情形不会发生。

为理解 $U$ 的结构，必须解开复杂波型。为此，波不仅须在相邻碰撞之间追踪，而须全局地从诞生追踪到消亡或永生。波通过如下约定被赋予全局身份：卷入碰撞的 $i$-波未必在该处终止，而一般作为从该波相互作用点出射的 $i$-波继续存在。应用上述规则时可能出现的歧义将在下面处理并解决。

伪激波由两个前沿碰撞产生、经简化 Riemann 解算器求解，如\cref{figure:ch14-不同族前沿碰撞}或\cref{figure:ch14-同族前沿碰撞}所示。另一方面，$i$-前沿既可在 $t=0$ 处由初始阶梯函数的某跳跃间断求解产生，也可在 $t>0$ 处由 $j$-前沿与 $k$-前沿（$j\neq i\neq k$）碰撞、经近似 Riemann 解算器求解而产生。

每个波在其整个生存期内携带一个数 $\mu$，标识其\emph{生成阶}，即早于其诞生的碰撞的最大数目。因而在 $t=0$ 处起源的前沿被赋予生成阶 $\mu=0$。任何其他新波必然由两个波碰撞产生，其生成阶分别为 $\mu_1$ 与 $\mu_2$，则被赋予生成阶 $\mu=\max\{\mu_1,\mu_2\}+1$。

如上所设，波穿过相互作用点时保持其生成阶。当经近似 Riemann 解算器求解的 $i$-稀疏前沿与 $j$-前沿碰撞中，出射 $i$-波扇含有两个 $i$-稀疏前沿时，可能出现歧义。此时这些前沿中较强者（强度为 $\delta$）被指定为入射 $i$-前沿的延续，而另一前沿（强度 $<\delta$）被视为新前沿，并按标准规则赋予更高生成阶。当同族两个前沿碰撞时也可能出现歧义，因为出射波扇至多含一个该族前沿。此情形下约定较低生成阶的前沿为存活者，另一个终止。若两前沿生成阶相同，则可任意指定其一为存活者。当然，如\cref{figure:ch14-同族前沿碰撞}所示，在非一般情形（一个为压缩激波、另一个为同族稀疏前沿且强度相同）下，两前沿可在碰撞时同时终止。伪激波也可如\cref{figure:ch14-伪激波与前沿碰撞}所示在非一般情形 $U_q=U_r$ 中通过与前沿碰撞而在有限时间内熄灭。

现引入下述概念，它们将建立与\cref{chapter:ch10-广义特征线}--\cref{chapter:ch12-两方程真非线性组}所用方法的联系。



对 $i=1,\cdots,n$，与 $U$ 相伴的\textbf{$i$-特征}是一条 Lipschitz 折线 $x=\xi(t)$，它在常状态 $\bar U$ 处以经典 $i$-特征速度 $\dot\xi=\lambda_i(\bar U)$ 穿过，但当撞上 $i$-前沿或其生成点时，便粘附于该前沿并在其整个生存期内跟随之。因而特别地，任一 $i$-前沿都是 $i$-特征。类似地，$(n+1)$-特征定义为斜率为 $\lambda_{n+1}$ 的直线。因而伪激波是 $(n+1)$-特征。

现考虑一条有向 Lipschitz 曲线，其图像 $\mathscr{C}$ 把上半平面分成``正''侧与``负''侧。称 $\mathscr{C}$ 为\textbf{非共振的}，若集合 $\{1,\cdots,n,n+1\}$ 可划分为三个两两不交（可能为空）的子集 $\mathscr{N}_-,\mathscr{N}_0,\mathscr{N}_+$，具有如下性质：$\mathscr{N}_-$ 与 $\mathscr{N}_+$ 各由至多 $n+1$ 个连续整数组成，而 $\mathscr{N}_0$ 至多含一个元素。对 $i\in\mathscr{N}_-$（或 $i\in\mathscr{N}_+$），任一撞上 $\mathscr{C}$ 的 $i$-特征从正侧穿到负侧（或从负侧穿到正侧）。另一方面，若 $i\in\mathscr{N}_0$，任一从正侧或负侧撞上 $\mathscr{C}$ 的 $i$-特征被 $\mathscr{C}$ 吸收，即 $\mathscr{C}$ 本身是一条 $i$-特征。

值得注意的非共振曲线例子包括：
\begin{enumerate}[(a)]
\item 任一 $i$-特征（特别地，任一 $i$-波），此时划分 $\mathscr{N}_-=\{1,\cdots,i-1\}$、$\mathscr{N}_0=\{i\}$、$\mathscr{N}_+=\{i+1,\cdots,n+1\}$。
\item 任一空间-like 曲线。设 $\lambda_1(U)<0<\lambda_{n+1}$，它们可表示为 Lipschitz 函数 $t=\hat t(x)$，a.e. 满足 $1/\lambda_1<\mathrm{d}\hat t/\mathrm{d}x<1/\lambda_{n+1}$。此情形下 $\mathscr{N}_+=\{1,\cdots,n+1\}$，而 $\mathscr{N}_-$ 与 $\mathscr{N}_0$ 均为空。
\end{enumerate}

上述内容的意义将在下一节中显现。

\end{document}
